%%%%%%%%%%%%%%%%%%%%%%%%%%%%%%%%%%%%%%%%%%%%%%%%%%%%%%%%%%%%
%% Lecture 09
%%%%%%%%%%%%%%%%%%%%%%%%%%%%%%%%%%%%%%%%%%%%%%%%%%%%%%%%%%%%

\lecture
{09}
{Monday, 7 September 2026}
{Introduction to Probability}
{Dr. Nishant Chandgotia}


%%%%%%%%%%%%%%%%%%%%%%%%%%%%%%%%%%%%%%%%%%%%%%%%%%%%%%%%%%%%
%% Continuation of Coupon Collector Problem
%%%%%%%%%%%%%%%%%%%%%%%%%%%%%%%%%%%%%%%%%%%%%%%%%%%%%%%%%%%%

\section{Continuation of Coupon Collector Problem}

We have already defined $X_n$ and
\(
    Y_k = \sum_{i=1}^{n} S_{i,k}.
\)

Now,
\[
\begin{aligned}
    \operatorname{Var}(Y_k)
    &=
    \sum_{i=1}^{n}\operatorname{Var}(S_{i,k})
    +
    \sum_{\substack{i,j=1\\i\neq j}}^{n}
    \operatorname{Cov}(S_{i,k},S_{j,k}).
\end{aligned}
\]

We have
\[
    \operatorname{Var}(S_{i,k})
    =
    \left(1-\frac{1}{n}\right)^k
    \left[
        1-\left(1-\frac{1}{n}\right)^k
    \right].
\]

Now,
\[
\begin{aligned}
    \operatorname{Cov}(S_{i,k},S_{j,k})
    &=
    \mathbb{E}(S_{i,k}S_{j,k})
    -
    \mathbb{E}(S_{i,k})\mathbb{E}(S_{j,k}) \\
    &=
    \left(1-\frac{2}{n}\right)^k
    -
    \left(1-\frac{1}{n}\right)^{2k}.
\end{aligned}
\]

Since
\[
    \left(1-x\right)^2
    \geq
    1-2x,
\]
we have
\[
    \left(1-\frac{1}{n}\right)^2
    \geq
    1-\frac{2}{n},
\]
and hence
\[
    \left(1-\frac{2}{n}\right)^k
    \leq
    \left(1-\frac{1}{n}\right)^{2k}.
\]

Therefore,
\[
    \operatorname{Cov}(S_{i,k},S_{j,k})
    \leq 0.
\]

Thus,
\[
\begin{aligned}
    \operatorname{Var}(Y_k)
    &\leq
    n
    \left(1-\frac{1}{n}\right)^k
    \left[
        1-\left(1-\frac{1}{n}\right)^k
    \right].
\end{aligned}
\]

Also,
\[
    \mathbb{E}(Y_k)
    =
    n\left(1-\frac{1}{n}\right)^k.
\]

Therefore,
\[
\begin{aligned}
    \frac{\mathbb{E}(Y_k)}
    {\sqrt{\operatorname{Var}(Y_k)}}
    &\geq
    \frac{
        n\left(1-\frac{1}{n}\right)^k
    }{
        \sqrt{
            n\left(1-\frac{1}{n}\right)^k
            \left[
                1-\left(1-\frac{1}{n}\right)^k
            \right]
        }
    } \\
    &\geq
    \sqrt{
        n\left(1-\frac{1}{n}\right)^k
    }.
\end{aligned}
\]

\begin{question}
{Coupon Collector Problem}
{coupon-collector-lower-bound}

We want to show that
\[
    \mathbb{P}
    \left(
        X_n
        <
        (1-\epsilon)n\log n
    \right)
    \longrightarrow 0
    \qquad\text{as }n\to\infty.
\]

Consider
\[
    k=(1-\epsilon)n\log n.
\]

Then
\[
\begin{aligned}
    \mathbb{P}(X_n<k)
    &=
    \mathbb{P}(Y_k=0).
\end{aligned}
\]

By Chebyshev's inequality,
\[
\begin{aligned}
    \mathbb{P}(Y_k=0)
    &\leq
    \mathbb{P}
    \left(
        \left|
        \frac{Y_k-\mathbb{E}(Y_k)}
        {\sqrt{\operatorname{Var}(Y_k)}}
        \right|
        \geq
        \frac{\mathbb{E}(Y_k)}
        {\sqrt{\operatorname{Var}(Y_k)}}
    \right) \\
    &\leq
    \frac{
        \operatorname{Var}(Y_k)
    }{
        \mathbb{E}(Y_k)^2
    }.
\end{aligned}
\]

More generally, Chebyshev gives
\[
    \mathbb{P}
    \left(
        \left|
        \frac{Y_k-\mathbb{E}(Y_k)}
        {\sqrt{\operatorname{Var}(Y_k)}}
        \right|
        \geq t
    \right)
    \leq
    \frac{1}{t^2}.
\]

We have
\[
\begin{aligned}
    \left|
        Y_k-\mathbb{E}(Y_k)
    \right|
    \geq
    t\sqrt{\operatorname{Var}(Y_k)}
    \quad\Longleftrightarrow\quad
    Y_k
    \notin
    \left(
        \mathbb{E}(Y_k)
        -
        t\sqrt{\operatorname{Var}(Y_k)},
        \mathbb{E}(Y_k)
        +
        t\sqrt{\operatorname{Var}(Y_k)}
    \right).
\end{aligned}
\]

Choose $t$ such that
\[
    0<
    \mathbb{E}(Y_k)
    -
    t\sqrt{\operatorname{Var}(Y_k)}.
\]

Thus,
\[
    t
    \leq
    \frac{\mathbb{E}(Y_k)}
    {\sqrt{\operatorname{Var}(Y_k)}}.
\]

If
\[
    Y_k=0,
\]
then
\[
\begin{aligned}
    \left|
        \frac{Y_k-\mathbb{E}(Y_k)}
        {\sqrt{\operatorname{Var}(Y_k)}}
    \right|
    &=
    \frac{\mathbb{E}(Y_k)}
    {\sqrt{\operatorname{Var}(Y_k)}} \\
    &\geq
    \sqrt{
        n\left(1-\frac{1}{n}\right)^k
    }.
\end{aligned}
\]

For
\[
    k=(1-\epsilon)n\log n,
\]
we have
\[
\begin{aligned}
    n\left(1-\frac{1}{n}\right)^k
    &\sim
    n e^{-(1-\epsilon)\log n} \\
    &=
    n^\epsilon.
\end{aligned}
\]

Hence,
\[
    \frac{\mathbb{E}(Y_k)}
    {\sqrt{\operatorname{Var}(Y_k)}}
    \geq
    \sqrt{n^\epsilon}
    =
    n^{\epsilon/2}.
\]

Therefore,
\[
    \mathbb{P}
    \left(
        Y_{(1-\epsilon)n\log n}=0
    \right)
    \longrightarrow 0
    \qquad\text{as }n\to\infty.
\]

\end{question}

\textbf{Todo:} Check and fill the corresponding graphs.


%%%%%%%%%%%%%%%%%%%%%%%%%%%%%%%%%%%%%%%%%%%%%%%%%%%%%%%%%%%%
%% Erdős--Rényi Graphs
%%%%%%%%%%%%%%%%%%%%%%%%%%%%%%%%%%%%%%%%%%%%%%%%%%%%%%%%%%%%

\section{Erdős--Rényi Graphs}

Let
\(
    n\in\mathbb{N},
    \qquad
    p\in(0,1).
\)
The random graph
\(
    G(n,p)
\)
is a random graph with $n$ vertices
\(
    \{1,2,\ldots,n\},
\)
where each edge is chosen independently with probability $p$.

\begin{question}
{Number of Copies of $K_r$}
{number-of-copies-kr}

How many copies of $K_r$ are there in $G(n,p)$, where $K_r$
denotes the complete graph on $r$ vertices?
\end{question}

Let
\(
    A\subseteq\{1,2,\ldots,n\}
\)
have cardinality $r$.

Define
\(
    X_A
    =
    \begin{cases}
        1, & \text{if $K_r$ appears on $A$},\\
        0, & \text{otherwise}.
    \end{cases}
\)

Let $X$ denote the number of copies of $K_r$. Then
\[
    X
    =
    \sum_{\substack{
        A\subseteq\{1,2,\ldots,n\}\\
        |A|=r
    }}
    X_A.
\]

Since all the edges of $K_r$ must be present,
\[
    \mathbb{E}(X_A)
    =
    p^{\binom{r}{2}}.
\]

Therefore,
\[
\begin{aligned}
    \mathbb{E}(X)
    &=
    \sum_{\substack{
        A\subseteq\{1,2,\ldots,n\}\\
        |A|=r
    }}
    \mathbb{E}(X_A) \\
    &=
    \binom{n}{r}
    p^{\binom{r}{2}}.
\end{aligned}
\]

Now,
\[
\begin{aligned}
    \operatorname{Var}(X)
    &=
    \sum_{\substack{
        A\subseteq\{1,2,\ldots,n\}\\
        |A|=r
    }}
    \operatorname{Var}(X_A) \\
    &\qquad+
    \sum_{\substack{
        A\neq B\\
        |A|=|B|=r
    }}
    \operatorname{Cov}(X_A,X_B).
\end{aligned}
\]

For the first term,
\[
\begin{aligned}
    \sum_{|A|=r}\operatorname{Var}(X_A)
    &=
    \binom{n}{r}
    p^{\binom{r}{2}}
    \left(
        1-p^{\binom{r}{2}}
    \right).
\end{aligned}
\]

If
\(
    |A\cap B|\leq 1,
\)
then the sets of edges involved in $X_A$ and $X_B$ are disjoint.

Hence,
\[
    \operatorname{Cov}(X_A,X_B)=0.
\]

If
\(
    |A\cap B|=t>1,
\)
then
\[
\begin{aligned}
    \operatorname{Cov}(X_A,X_B)
    &=
    \mathbb{E}(X_AX_B)
    -
    \mathbb{E}(X_A)\mathbb{E}(X_B) \\
    &=
    p^{
        2\binom{r}{2}-\binom{t}{2}
    }
    -
    p^{2\binom{r}{2}}.
\end{aligned}
\]

Thus,
\[
\begin{aligned}
    \sum_{\substack{
        A\neq B,\,
        |A\cap B|=t\\
        |A|=|B|=r
    }}
    \operatorname{Cov}(X_A,X_B)
    &=
    \binom{n}{r}
    \binom{n-r}{r-t}
    \left[
        p^{
            2\binom{r}{2}-\binom{t}{2}
        }
        -
        p^{2\binom{r}{2}}
    \right].
\end{aligned}
\]

Consequently,
\[
\begin{aligned}
    \operatorname{Var}(X)
    &=
    \binom{n}{r}
    p^{\binom{r}{2}}
    \left(
        1-p^{\binom{r}{2}}
    \right) \\
    &\qquad+
    \sum_{t=2}^{r-1}
    \binom{n}{r}
    \binom{n-r}{r-t}
    \left[
        p^{
            2\binom{r}{2}-\binom{t}{2}
        }
        -
        p^{2\binom{r}{2}}
    \right].
\end{aligned}
\]

For some constants $c_1,c_2>0$,
\[
    \mathbb{E}(X)\sim c_1 n^r,
    \qquad
    \operatorname{Var}(X)\sim c_2 n^{2r-2}.
\]

By Chebyshev's inequality, we have
\[
    \mathbb{P}
    \left(
        \left|
            \frac{X-\mathbb{E}(X)}
            {n^{r-1}f(n)}
        \right|
        >\epsilon
    \right)
    \longrightarrow 0
    \qquad\text{as }n\to\infty,
\]
where $f(n)$ is any function such that
\(
    f(n)\longrightarrow\infty
    \qquad\text{as }n\to\infty.
\)

In fact,
\[
\begin{aligned}
    \mathbb{P}
    \left(
        |X-\mathbb{E}(X)|
        >
        \epsilon\mathbb{E}(X)
    \right)
    &\leq
    \frac{
        \operatorname{Var}(X)
    }{
        \epsilon^2\mathbb{E}(X)^2
    } \\
    &\sim
    \frac{1}{n^2\epsilon^2}
    \longrightarrow 0.
\end{aligned}
\]

%%%%%%%%%%%%%%%%%%%%%%%%%%%%%%%%%%%%%%%%%%%%%%%%%%%%%%%%%%%%
%% Full Strength of Chebyshev
%%%%%%%%%%%%%%%%%%%%%%%%%%%%%%%%%%%%%%%%%%%%%%%%%%%%%%%%%%%%

\subsection{Full Strength of Chebyshev}

Let
\(
    X_1,X_2,\ldots,X_n
\)
be random variables and define
\(
    S_n=\sum_{i=1}^{n}X_i.
\)

Suppose
\(
    \mathbb{E}(X_i)=\mu,
\)
\(
    \operatorname{Var}(X_i)\leq C
    \qquad\text{for all }i,
\)
and
\[
    \operatorname{Cov}(X_i,X_j)=0
    \qquad\text{for all }i<j.
\]

Thus, the random variables are uncorrelated.

Then
\(
    \frac{S_n}{n}
    \longrightarrow
    \mu
\)
in $L^2$ and in probability.

This is a version of the weak law of large numbers.

\textbf{Proof.}

We have
\[
\begin{aligned}
    \mathbb{E}
    \left[
        \left|
            \frac{S_n}{n}-\mu
        \right|^2
    \right]
    &=
    \operatorname{Var}
    \left(
        \frac{S_n}{n}
    \right) \\
    &=
    \frac{1}{n^2}
    \operatorname{Var}(S_n).
\end{aligned}
\]

Since the random variables are uncorrelated,
\[
\begin{aligned}
    \operatorname{Var}(S_n)
    &=
    \sum_{i=1}^{n}\operatorname{Var}(X_i) \\
    &\leq
    Cn.
\end{aligned}
\]

Therefore,
\[
    \mathbb{E}
    \left[
        \left|
            \frac{S_n}{n}-\mu
        \right|^2
    \right]
    \leq
    \frac{Cn}{n^2}
    =
    \frac{C}{n}
    \longrightarrow 0.
\]

Hence,
\(
    \frac{S_n}{n}
    \longrightarrow
    \mu
    \qquad\text{in }L^2,
\)
and therefore
\(
    \frac{S_n}{n}
    \longrightarrow
    \mu
    \qquad\text{in probability}.
\)

\textbf{But further:}

By Chebyshev's inequality,
\[
\begin{aligned}
    \mathbb{P}
    \left(
        \left|
            \frac{S_n}{n}-\mu
        \right|
        >
        \epsilon\mu
    \right)
    &\leq
    \frac{
        \operatorname{Var}(S_n)
    }{
        (n\epsilon\mu)^2
    } \\
    &\leq
    \frac{nC}
    {n^2\epsilon^2\mu^2} \\
    &=
    \frac{C}
    {n\epsilon^2\mu^2}
    \longrightarrow 0.
\end{aligned}
\]

\begin{question}
{Exponential Decay}
{exponential-decay}

How do we get exponential decay?
\end{question}


\lectureend
